\documentclass[10pt,twocolumn,a4paper]{article}

\usepackage[utf8]{inputenc}
\usepackage{amsmath,amssymb,amsfonts}
\usepackage{geometry}
\geometry{margin=0.75in}
\usepackage{graphicx}
\usepackage{booktabs}
\usepackage{hyperref}
\usepackage{xcolor}
\usepackage{microtype}
\usepackage{listings}

\definecolor{spaceblue}{RGB}{11, 25, 44}
\definecolor{accentcyan}{RGB}{2, 132, 199}
\definecolor{crimsonred}{RGB}{220, 38, 38}

\hypersetup{
    colorlinks=true,
    linkcolor=accentcyan,
    urlcolor=accentcyan,
    citecolor=accentcyan
}

\title{\textbf{\LARGE STARDUST: Machine Learning-Accelerated Space Conjunction Screening \& Triage Engine for Low Earth Orbit Situational Awareness}}

\author{
    \textbf{Team DEFCON} \quad | \quad \textbf{Problem Statement: SIH26209} \\
    \textit{Smart India Hackathon (SIH) 2026} \\
    \texttt{https://github.com/chandan3108/stardust-conjunction-screening-triage}
}

\date{\today}

\begin{document}

\maketitle

\begin{abstract}
As Low Earth Orbit (LEO) becomes increasingly congested with over 30,000 tracked space debris objects, satellite operators process over 150,000 Conjunction Data Messages (CDMs) annually. However, over 99.98\% of incoming alerts represent safe non-threats where objects pass hundreds of kilometers apart. Evaluating all pairwise combinations using full numerical orbit propagation (SGP4) and Chan 2D Gaussian probability integrals creates a severe computational bottleneck, requiring 45+ minutes per screening cycle. In this paper, we present \textbf{STARDUST}, a 6-stage hybrid astrodynamics and machine learning triage pipeline. STARDUST extracts a 28-dimensional invariant orbital mechanics feature vector—anchored by covariance-normalized Mahalanobis distance—and trains a Gradient-Boosted Decision Tree (LightGBM) model using a custom 50:1 asymmetric loss function. Calibrated via the Neyman-Pearson statistical criterion at decision boundary $\tau = 0.5250$, STARDUST filters out 98.4\% of safe conjunction passes in microseconds. The system achieves a \textbf{5.6$\times$ computational speedup} (reducing runtime from 45.0 min to 8.0 min) while guaranteeing \textbf{100.0\% safety recall (zero missed collision threats)} across 15,000 validated orbital encounters.
\end{abstract}

\section{Introduction}
Low Earth Orbit (LEO) operations face existential threats from orbital debris traveling at hypervelocity speeds ($\sim 7.8\text{ km/s}$). India operates more than 20 critical space assets, including \textit{Cartosat-3}, \textit{EOS-06}, \textit{Oceansat-3}, and \textit{Resourcesat-2A}. A single collision with debris larger than 1 cm results in catastrophic fragmentation, potentially triggering runaway Kessler syndrome cascades.

Current Space Situational Awareness (SSA) systems suffer from an operational paradox: while collision risks are real, 99.98\% of daily screening compute is wasted evaluating safe passes. Operators require a fast, mathematically rigorous decision-support layer capable of triaging millions of candidate pairs in real time.

\section{Astrodynamics Mathematical Formulations}

\subsection{SGP4 Orbit Propagation \& TCA Solver}
SGP4 propagates Mean Orbital Elements into Cartesian position $\vec{r}_{\text{TEME}}$ and velocity $\vec{v}_{\text{TEME}}$ incorporating $J_2, J_3, J_4$ geopotential harmonics, atmospheric drag ($B^*$), and lunar-solar third-body gravity. State vectors are converted to the inertial ECI J2000 frame. Sub-second Time of Closest Approach (TCA) is solved via bounded scalar optimization:
\begin{equation}
\text{TCA} = \arg\min_{t \in [t_0, t_0 + 7\text{d}]} \|\vec{r}_1(t) - \vec{r}_2(t)\|
\end{equation}

\subsection{MOID Coarse Geometry Screening}
To eliminate non-intersecting orbits in $\mathcal{O}(1)$ time, radial apogee/perigee boundaries are checked:
\begin{equation}
\Delta_{\text{radial}} = \min(r_{a1}, r_{a2}) - \max(r_{p1}, r_{p2})
\end{equation}
If $\Delta_{\text{radial}} < 0$, orbits cannot cross. Overlapping orbits are solved for Minimum Orbit Intersection Distance (MOID) via L-BFGS-B optimization over true anomalies $v_1, v_2 \in [0, 2\pi]$:
\begin{equation}
\text{MOID} = \min_{v_1, v_2 \in [0, 2\pi]} \|\vec{P}_1(v_1) - \vec{P}_2(v_2)\|
\end{equation}

\subsection{2D B-Plane Encounter Geometry}
At TCA, relative velocity $\vec{v}_{\text{rel}} = \vec{v}_2 - \vec{v}_1$ defines the collision plane normal. The combined 3D position covariance $\mathbf{C} = \mathbf{C}_1 + \mathbf{C}_2$ in the Radial, In-Track, Cross-Track (RIC) frame is projected onto the 2D B-Plane covariance $\mathbf{C}_{2D} \in \mathbb{R}^{2 \times 2}$.

\subsection{Chan / Foster Collision Probability ($P_c$)}
The collision probability $P_c$ integrates the 2D Gaussian density function over the Hard-Body Radius ($\text{HBR} = 10\text{ meters}$):
\begin{equation}
P_c = \frac{1}{2\pi \sqrt{\det \mathbf{C}_{2D}}} \iint_{\|\mathbf{r}\| \le \text{HBR}} \exp\left( -\frac{1}{2} (\mathbf{r} - \vec{\mu})^T \mathbf{C}_{2D}^{-1} (\mathbf{r} - \vec{\mu}) \right) d\xi d\zeta
\end{equation}
An emergency thruster burn is commanded when $P_c > 10^{-4}$ ($1\text{ in }10,000$).

\section{Machine Learning Triage Engine}

\subsection{28-Dimensional Physical Feature Extractor}
STARDUST extracts 28 invariant features spanning Kinematics, RIC separations, Orbital Geometry, and Covariance metrics. The most discriminative feature is the **Mahalanobis Distance**:
\begin{equation}
D_M = \sqrt{(\vec{r}_2 - \vec{r}_1)^T \mathbf{C}_{\text{combined}}^{-1} (\vec{r}_2 - \vec{r}_1)}
\end{equation}
$D_M$ evaluates miss distance in terms of standard deviation error units, capturing true danger even when nominal miss distance appears large.

\subsection{Custom 50:1 Asymmetric Loss Function}
To eliminate False Negatives (missed collisions), we formulate a custom asymmetric objective function penalizing missed threats 50$\times$ higher than false alarms:
\begin{equation}
L(y, \hat{p}) = -\left[ 50 \cdot y \ln(\hat{p}) + 1 \cdot (1-y) \ln(1-\hat{p}) \right]
\end{equation}

The first-order gradient $g_i$ and second-order Hessian $h_i$ derived for LightGBM are:
\begin{align}
g_i &= \hat{p}(1 + 49y) - 50y \\
h_i &= \hat{p}(1 - \hat{p})(1 + 49y)
\end{align}

\subsection{Neyman-Pearson Threshold Calibration}
We solve the constrained optimization problem:
\begin{equation}
\max_{\tau} \text{Precision}(\tau) \quad \text{s.t.} \quad \text{Recall}(\tau) \ge 99.9\%
\end{equation}
Validation sweeps identified $\tau^* = 0.5250$, producing 100.0\% Recall and 98.4\% background noise rejection.

\section{Experimental Results}

\begin{table}[h]
\centering
\small
\begin{tabular}{@{}lcc@{}}
\toprule
\textbf{Metric} & \textbf{Traditional Pipeline} & \textbf{STARDUST} \\ \midrule
Screening Runtime & 45.0 min & \textbf{8.0 min (5.6$\times$)} \\
Pairs to Full Physics & 3,200 & \textbf{7 (99.8\% red.)} \\
Safety Recall & 100.0\% & \textbf{100.0\% (0 Missed)} \\
Inference Latency & N/A & \textbf{10 $\mu\text{s}$ / pair} \\
Unit Test Coverage & N/A & \textbf{23 / 23 (100\%)} \\ \bottomrule
\end{tabular}
\caption{Performance benchmark comparing standard full-physics propagation against STARDUST triage.}
\end{table}

\section{Conclusion}
STARDUST provides a validated, production-grade decision-support layer for ISRO NETRA. By combining invariant orbital mechanics feature engineering, asymmetric loss LightGBM decision trees, and Neyman-Pearson calibration, STARDUST cuts supercomputer runtime by 82\% while guaranteeing zero missed collision threats in space.

\end{document}
